\bookchapter{Глава 40. Инженерия требований — превратите идеи в инструкции, выполнимые для ИИ}

\chapterintro{Если инженерия контекста решает задачу «дать ИИ понять фон», то инженерия требований решает «дать ИИ знать, что именно нужно сделать сейчас». Это этап, на котором разница между людьми видна сильнее всего: с теми же инструментами и той же моделью один делает восхитительные проекты, а другой — посредственные игрушки. Дело обычно не в технических навыках, а в умении выразить требования.}

\booksection{От смутной идеи к структурированному требованию}

Каждый проект начинается со смутной идеи — «хочу сделать приложение для учёта расходов», «хочу личный сайт», «хочу инструмент-список дел». Сами по себе эти идеи нормальные, но если скормить их ИИ напрямую, качество результата целиком зависит от удачи. Цель инженерии требований — превратить смутную идею в структурированное требование, которое ИИ сможет выполнить точно.

Хорошая новость: структурированные требования не требуют многотомного PRD. Достаточно ответить на пять вопросов — мы называем их «пятью элементами требования»:

\begin{center}
\begin{tabularx}{\linewidth}{|X|X|X|}
\hline
\textbf{Элемент} & \textbf{Вопрос} & \textbf{Пример: приложение для учёта расходов} \tabularnewline
\hline
Who & Для кого? & Личное пользование, учёт повседневных доходов и расходов \tabularnewline
What & Что делает? & Приложение, которое записывает каждую трату, сводит по месяцам и визуализирует данные \tabularnewline
Why & Зачем делать? & Хочу понять структуру своих трат и найти лишние расходы \tabularnewline
How & Как примерно делать? & Фронтенд на React, данные храним локально в браузере, бэкенд и вход пока не делаем \tabularnewline
Constraint & Какие ограничения? & Один разработчик, готово за 2 недели, деплой в продакшен не нужен \tabularnewline
\hline
\end{tabularx}
\end{center}

\begin{center}\small\color{gray}Пять элементов требования и пример\end{center}

Эти пять элементов не нужно оформлять в официальный документ — достаточно быстро прокрутить их в голове. Главное — выработать привычку «сначала подумать, потом говорить». Вы заметите: когда вы способны ясно ответить на пять вопросов, описание требований само собой становится структурированным.

\booksection{Искусство послойной декомпозиции требований}

Для крупных проектов одного описания требований мало. Большое требование нужно разбить на маленькие подзадачи, каждая из которых проверяется независимо. Этот процесс называется «декомпозицией требований»; его главный принцип: каждая подзадача должна быть выполнима, проверяема и сдаваема отдельно.

На примере «приложения для личного учёта расходов» разумная декомпозиция такая:

\begin{enumerate}
  \item Реализовать базовую запись доходов и расходов (добавлять, просматривать, удалять записи);
  \item Добавить категории (помечать каждую запись: еда, транспорт и т. д.);
  \item Реализовать месячные сводки (статистика доходов и расходов по месяцам);
  \item Добавить визуализацию данных (показывать структуру трат графиками).
\end{enumerate}

Каждый шаг — функциональное приращение, которое можно запустить и проверить отдельно.

Золотое правило декомпозиции требований: после каждого шага проект должен оставаться «работающей вещью» — пусть даже грубой. Никогда не оставляйте проект в состоянии «всё сделано наполовину, но ничего не работает».

У такой декомпозиции два плюса. Во-первых, вы всегда видите прогресс: с каждым шагом у приложения появляется новая функция — эта положительная обратная связь критична для долгих проектов (ощущение успеха нужно всем). Во-вторых, каждая задача ИИ остаётся достаточно маленькой, что снижает риск «непонимания со стороны ИИ» — чем меньше задача, тем точнее её понимает ИИ и тем выше качество результата.

\booksection{Шаблон PRD и практика написания}

Для крупных проектов (например, полноценного приложения с фронтендом, бэкендом и несколькими модулями) пяти элементов требований может не хватить. Тогда нужен более официальный документ требований проекта (PRD). Частое недопонимание: PRD пишется не для ИИ — он пишется для вас самих, чтобы прояснить мысли, и уже тогда вы сможете дать ИИ более точные инструкции.

Хороший PRD не обязан быть длинным, но должен содержать шесть частей: обзор проекта (одним абзацем — что делаем), список функций (что представляет каждая функция), техническое решение (какой стек, какая архитектура в общих чертах), модель данных (ключевые структуры данных), описание интерфейса (как выглядят главные экраны), план разработки (за сколько шагов сделаем).

\begin{lstlisting}
# PRD приложения личных финансов
## О проекте
Лёгкое приложение личных финансов: записи доходов/расходов, категории, месячные итоги и визуализация.
## Список функций
1. Запись доходов/расходов: сумма, тип (доход/расход), категория, заметка, дата
2. Управление категориями: предустановки — еда, транспорт, покупки, развлечения, другое
3. Итоги по месяцам: общий доход, общий расход, остаток
4. Визуализация расходов: круговая диаграмма долей категорий
5. Экспорт данных: выгрузка в CSV
## Техническое решение
- Фронтенд: React + TypeScript + Tailwind CSS
- Хранение данных: localStorage (бэкенд не нужен)
- Диаграммы: Chart.js
- Сборка: Vite
## Модель данных
Record {
  id: string
  amount: number
  type: 'income' | 'expense'
  category: string
  note: string
  date: string  // YYYY-MM-DD
}
## План разработки
Step 1: базовые функции записей (CRUD)
Step 2: управление категориями
Step 3: итоги по месяцам
Step 4: визуализация данных
Step 5: экспорт данных
\end{lstlisting}

Обратите внимание на стиль этого PRD: краткие списки и структурированный формат вместо многословия. У каждой функции — чёткое определение, модель данных описана простыми типами. Такой PRD вам самим легко читать, а когда вы передадите его ИИ, тот точно поймёт каждый пункт требований.

\booksection{Практика: полный прогон от идеи до решения}

Рассуждать на пальцах скучно — пройдём весь процесс инженерии требований на цельном примере. Допустим, ваша идея: «хочу сделать трекер учёбы, который фиксирует, что я каждый день выучил».

Шаг первый: пробегаем пять элементов требования.

\begin{itemize}
  \item Who: личное пользование.
  \item What: инструмент, который записывает, что выучено за день, отслеживает время учёбы и помогает пересматривать прогресс.
  \item Why: хочу выработать привычку учиться без перерывов и видеть свои накопления.
  \item How: сначала веб-версия, данные хранятся локально в браузере.
  \item Constraint: разработка с нуля, ключевые функции — за 1 неделю.
\end{itemize}

Шаг второй: разбираем требования на части.

\begin{enumerate}
  \item Step 1: реализовать функцию «записи учебных пунктов» (дата, содержание, длительность);
  \item Step 2: добавить «статистику времени учёбы» (за сегодня, суммарно за неделю);
  \item Step 3: реализовать вид «учебного календаря» (как график вкладов GitHub);
  \item Step 4: добавить функцию «повторения» (случайно показывать записи учёбы за какой-нибудь прошлый день, помогая вспоминать).
\end{enumerate}

Шаг третий: напишите короткий PRD. Длинный не нужен — сама декомпозиция выше уже составляет тело PRD. Оформите его файлом .md в корне проекта как общий справочник для вас и ИИ.

Шаг четвёртый: начинайте разработку первого шага. Инструкция для ИИ может быть такой:

\begin{lstlisting}
Нужна первая функция трекера обучения — запись учебных пунктов.
Пользователь может добавить запись: дата, что изучал, длительность в минутах.
Данные хранятся в localStorage.
Фронтенд — React + TypeScript.
Дайте полную реализацию.
\end{lstlisting}

Обратите внимание на особенности этой инструкции: она описывает только цель текущего шага (функцию записи), не примешивая другие функции; она задаёт технические ограничения (React, TypeScript, localStorage); она даёт достаточно информации, чтобы ИИ понял «что делать» и «зачем», но конкретную реализацию оставляет ИИ. Это и есть высококачественная инструкция для ИИ-программирования, обработанная инженерией требований.

\begin{tcbpractice}{Практика: составьте карточку пяти элементов для своей «заветной идеи»}
Выберите маленький проект, который вынашиваете давно, и за 5 минут заполните карточку: Who / What / Why / How / Constraint — по одной фразе на каждый. Затем разбейте его на 3–5 шагов, после каждого из которых «уже можно пользоваться». Вы увидите: стоит записать идею — и она превращается из «грёз наяву» в «список дел».

\end{tcbpractice}

\begin{tcbsummary}{Краткие итоги главы}
\begin{itemize}
  \item Пять элементов требования: Who / What / Why / How / Constraint — сначала подумайте, потом говорите
  \item Золотое правило декомпозиции: после каждого шага проект остаётся «работающей вещью»
  \item PRD пишется для себя: обзор, список функций, техрешение, модель данных, интерфейс, план
  \item Хорошая инструкция = только текущий шаг + ясные ограничения + говорим «что делать», а не «как делать»
  \item Шаблон полного прогона: пять элементов декомпозиция PRD инструкция, всего четыре шага
\end{itemize}

\end{tcbsummary}
